\section{Discussion}
  \label{sec:discussion}

  \paragraph{The conceptual content of the Gram identity.}
    The proof of Theorem~\ref{thm:gram-rigidity} rests entirely on two character-theoretic
    objects: $M = \sum \vec{u}_s\vec{u}_s^T$ and $(G_{st})$.
    Both are computable from the character table of $G$ and the representation $\rho$,
    with no geometric input.
    The Gram identity $G_{st} = -\tfrac{1}{2}\chi_\rho(st)$ is the key step:
    it shows that the pairwise \emph{angles} between neutral axes are encoded in the
    \emph{characters of pairwise products}.
    Geometry emerges from spectral arithmetic.

  \paragraph{Character-theoretic reconstruction of the Klein list.}
    The Klein classification of finite subgroups of $\SU(2)$ is classically proved
    by polyhedral geometry or by Dynkin diagram methods.
    The present paper gives a \emph{character-theoretic} proof of the same classification,
    restricted to the neutral sector.
    The three regimes of Theorem~\ref{thm:gram-rigidity} correspond to:
    \begin{itemize}
      \item an \emph{anisotropy} discriminant ($M \propto I$ or not), detecting the
      presence of a preferred rotation axis;
      \item an \emph{orthogonality} discriminant ($G_{st} = 0$ or not), detecting
      the purely quaternionic geometry;
      \item a \emph{non-triviality} discriminant ($G_{st} \in \{0,\pm 1\}$ or not),
      detecting the polyhedral (octahedral or icosahedral) geometry.
    \end{itemize}
    Each discriminant is a statement about eigenvalues of products of group elements,
    hence a purely spectral condition.

  \paragraph{Connection to the Admissibility Programme.}
    In the \textsc{Cosmochrony} framework, the generator set $S$ models the flux-carrying
    modes of the substrate operator $\mathcal{L}_\chi$, and the neutrality condition
    $\chi_\rho(s) = 0$ encodes compatibility with the bounded-flux constraint $\|c_{\mathrm{BI}}\| \le 1$
    (Steps~1--2).
    The present result shows that this constraint is powerful enough to recover the
    full polyhedral structure of $\SU(2)$: the admissible groups~\eqref{eq:admissible-list}
    are exactly those whose character geometry is consistent with isotropy or
    orthogonality of the axis configuration.

  \paragraph{Structural role of $2T$.}
    The exclusion of $2T$ is not an accident of numerics but a structural consequence of
    the coset decomposition: the tetrahedral group's neutral sector is entirely contained
    in $Q_8$, making it impossible for any neutral generating set to witness the full $A_4$
    symmetry.
    This asymmetry between $2T$ and the binary octahedral/icosahedral groups $2O$, $2I$
    reflects a deeper distinction: for $2O$ and $2I$, the full set of traceless elements
    \emph{is} a neutral generating set for the group, while for $2T$ it is not.

  \paragraph{Scale stability of the classification.}
    The classification established in Theorem~7.1 is derived under the Born--Infeld admissibility
    bound $A_n \leq c_{\mathrm{BI}}/\sqrt{\lambda_n}$, which defines the neutrality window in terms of the
    primitive scale $c_{\mathrm{BI}}$.
    The companion note~\cite{NoScaleL1} establishes that this window is structurally fixed:
    no independent dimensional parameter can be introduced into the admissibility constraints
    without breaking at least one of BFS-consistent spectral scaling, structural non-injectivity,
    or bounded admissibility closure.
    The group classification of Theorem~7.1 is therefore not contingent on a particular choice
    of admissibility scale, but is universally rigid within the Cosmochrony framework.
